% ============================================================================
% DRAFT v2 -- Section 3: The pipeline under study
% Reports what the pipeline does, not what it is worth.
% ============================================================================

\section{The Pipeline Under Study}
\label{sec:method}

This section describes the pipeline whose two design axes are examined in
Sections~\ref{sec:diagnosis} and~\ref{sec:leverage}. It is a three-stage
knowledge-distillation pipeline built on reverse
distillation~\cite{deng2022reverse}: a teacher is pre-trained on normal images,
an architecture is searched using a distillation signal, and the winning
architecture is trained as a reverse-distillation detector.
Figure~\ref{fig:pipeline} sketches the three stages and marks the two axes this
paper measures. We describe the pipeline precisely, because several of the
findings that follow are properties of specific implementation decisions rather
than of the method in the abstract. Stage~1 pre-trains a teacher encoder on normal
training images with a contrastive objective. Stage~2 searches an architecture for
the student, scoring each candidate by how well it distils from the frozen teacher
on normal data using a small number of training epochs. Stage~3 trains the
selected architecture --- student encoder plus decoder --- as a reverse-distillation
detector. The search runs on a single category and its result is applied to all
categories of the dataset.

\subsection{Stage 1: Teacher Pre-training}
\label{subsec:teacher}
% ============================================================================
% Figure 1 -- the pipeline under study, and the scoring policy.
% Vector (TikZ); no raster resolution limit applies.
% Included with \input from Section 3 (method).
% ============================================================================

\begin{figure*}[t]
\centering
\resizebox{\textwidth}{!}{%
\begin{tikzpicture}[
  font=\small,
  node distance=4mm,
  box/.style={draw, rounded corners=1pt, align=center, inner sep=3pt,
              minimum height=9mm, line width=0.5pt},
  stage/.style={box, fill=black!4, minimum width=20mm},
  data/.style={box, fill=blue!6, minimum width=17mm},
  sel/.style={box, fill=orange!14, minimum width=20mm},
  arr/.style={->, >=latex, line width=0.5pt},
  lab/.style={font=\scriptsize, inner sep=1pt},
]

%% ---------- Stage 1 ----------
\node[data] (imgs) {normal\\images};
\node[stage, right=6mm of imgs] (teacher) {contrastive\\pre-training};
\node[data, right=6mm of teacher] (frozen) {frozen teacher\\$t_1,t_2,t_3$};

\draw[arr] (imgs) -- (teacher);
\draw[arr] (teacher) -- (frozen);
\node[lab, above=0.5mm of teacher] {Stage 1};
\node[lab, below=0.5mm of teacher] {100 epochs};

%% ---------- Stage 2 ----------
\node[sel, below=11mm of teacher] (space) {42 configurations\\14 backbones $\times$ 3 dims};
\node[sel, right=6mm of space] (proxy) {KD proxy\\$s(a)=\overline{\mathrm{kd}+0.3\,\mathrm{cons}}$};
\node[sel, right=6mm of proxy] (argmin) {$\arg\min s(a)$\\single sample};

\draw[arr] (space) -- (proxy);
\draw[arr] (proxy) -- (argmin);
\node[lab, below=0.5mm of proxy] {3 epochs, wood};
\node[lab, above=0.5mm of space, xshift=6mm] {Stage 2};

%% ---------- Stage 3 ----------
\node[stage, below=11mm of proxy] (rd) {train student $+$ decoder\\500 epochs};
\node[data, right=6mm of rd, minimum width=23mm] (maps) {discrepancy maps\\$m_1,m_2,m_3$};

\draw[arr] (rd) -- (maps);
\node[lab, above=0.5mm of rd, xshift=8mm] {Stage 3};

%% ---------- the scoring policy ----------
\node[data, right=6mm of maps] (pool) {$S(x)=\rho\!\left(\sum_k w_k m_k\right)$\\[1pt]
     $\rho\in\{\text{mean},\text{top-}k,\max\}$};

\draw[arr] (maps) -- (pool);

%% ---------- arrows between stages ----------
\draw[arr] (argmin.south) |- ([yshift=2mm]rd.east) -- (rd.east);
\draw[arr] (frozen.east) -| ([xshift=3mm]space.north) -- (space.north);

%% ---------- annotation: what the paper measures ----------
\node[lab, align=left, text width=44mm, anchor=north west]
  at ([xshift=-2mm,yshift=-4mm]space.south west) {%
  \textcolor{orange!80!black}{Sections 4.1--4.3:} the searched dimension does not
  reach the detector; the proxy does not separate architectures; the axis carries
  $0.54\%$ of the variance.};

\node[lab, align=left, text width=40mm, anchor=north west]
  at ([xshift=2mm,yshift=-4mm]pool.south west) {%
  \textcolor{green!45!black}{Section 5:} the reduction $(\mathbf{w},\rho)$ is the
  axis that moves accuracy, and it can be chosen without labels.};

\end{tikzpicture}
}
\caption{The pipeline under study. A teacher is pre-trained on normal images
(Stage~1); an architecture is selected from 42 configurations by minimising a
short-training distillation proxy on one category (Stage~2); the selected
architecture is then trained as a reverse-distillation detector (Stage~3) and each
test image is scored by reducing its multi-scale discrepancy maps with a policy
$(\mathbf{w},\rho)$. This paper measures the two axes in bold: the architecture
that Stage~2 selects, and the reduction that Stage~3 applies.}
\label{fig:pipeline}
\end{figure*}


The teacher is a ResNet-18~\cite{he2016deep} with a 128-dimensional projection
head, trained for 100 epochs on the normal training images of the search category
with a contrastive objective, and frozen thereafter. Its role is to define the
feature space that both the search signal and the final detector compare against.
At test time it contributes intermediate features at three resolutions: the
outputs of \texttt{layer1}, \texttt{layer2} and \texttt{layer3}, of shape
$64\times56\times56$, $128\times28\times28$ and $256\times14\times14$ for an input
resized to $224\times224$.

\subsection{Stage 2: Architecture Search}
\label{subsec:nas}

The search space is the Cartesian product of 14 backbones and three feature
dimensions, giving 42 configurations. The backbones span convolutional and
attention-based families from 2015 to 2023, including
ResNets~\cite{he2016deep}, EfficientNet variants, MobileNetV3 and MobileOne, ViT
and DeiT variants, Swin-Tiny, and the ConvNeXt and ConvNeXt~V2
families~\cite{woo2023convnextv2}; Section~\ref{sec:positioning} lists them with
their deployed sizes. The three feature dimensions are 128, 256 and 512.

\subsubsection*{The evaluation signal}
Each candidate is evaluated by building a student with the candidate's backbone
and projection dimension, training it for a small number of epochs on the normal
training images of the search category, and then measuring a composite
distillation loss on held-out normal images:
%
\begin{equation}
\label{eq:proxy}
s(a) \;=\; \frac{1}{|\mathcal{V}|}\sum_{x \in \mathcal{V}}
\Bigl[\, \mathcal{L}_{\mathrm{kd}}(x) \;+\; 0.3\,\mathcal{L}_{\mathrm{cons}}(x) \Bigr],
\end{equation}
%
where $\mathcal{L}_{\mathrm{kd}}$ is the mean-squared error between the student's
projection and the teacher's, and $\mathcal{L}_{\mathrm{cons}}$ is the
mean-squared error between the student's projections of two independently
augmented views of the same image. Lower is better. Every candidate is scored once;
there is no repetition and therefore no estimate of the signal's own variance.
Two details of this computation are examined in
Section~\ref{subsec:proxy_stability} and are stated here for completeness. A
validation batch contributes to the sum only if its loss is finite and below a
fixed threshold; otherwise it is skipped. And the average is formed as
$\text{total} / \max(\text{count}, 1)$, so if every batch of a trial is skipped the
signal returns exactly zero. Because the backbones produce features of different
dimensionality while the teacher projects into a 128-dimensional space, a
candidate whose projection dimension differs from 128 is additionally given a
learned linear layer that maps the student's output into the teacher's space
before the losses are computed; this alignment layer exists only during search.

The search draws candidates uniformly at random from the 42 cells and returns the
one with the smallest score,
%
\begin{equation}
\label{eq:argmin}
a^\star \;=\; \arg\min_{a \in \mathcal{A}_{\text{sampled}}} \; s(a),
\end{equation}
%
and the number of draws is a protocol parameter (the recorded runs used 100, of
which a variable number complete successfully).

\subsection{Stage 3: Reverse-Distillation Detection}
\label{subsec:rd_detection}

The selected configuration is re-instantiated as a detector. A student encoder
takes the backbone's final feature map, adaptively pools it to $14\times14$, and
projects it through a $1\times1$ convolution with batch normalisation and ReLU to
a bottleneck of width $B$. A decoder then reconstructs the teacher's features in
reverse order: a first block produces the $14\times14$ stage at 256 channels, the
output is upsampled and passed through a second block producing the $28\times28$
stage at 128 channels, and upsampled again through a third block producing the
$56\times56$ stage at 64 channels. We note for later use that the bottleneck width
is fixed at $B = 256$ in the deployment path, independently of the feature
dimension that the search may have selected;
Section~\ref{subsec:inert} shows that this makes the searched dimension
unobservable in the deployed detector. The detector is trained on the normal
training images of each category separately, with AdamW at a learning rate of
$10^{-3}$ and weight decay $10^{-4}$, a cosine schedule, and random horizontal
flipping with probability $0.5$; the loss is the sum of mean-squared errors
between the student's reconstructions and the teacher's features at the three
scales.

\subsubsection*{Anomaly scoring}
\label{subsubsec:scoring}
Given a test image $x$, the teacher produces features $t_1, t_2, t_3$ and the
student produces reconstructions $s_1, s_2, s_3$. For each scale $k$ we form a
discrepancy map
%
\begin{equation}
\label{eq:diffmap}
m_k(x) \;=\; \bigl\| s_k(x) - t_k(x) \bigr\|_2 ,
\end{equation}
%
taking the $L_2$ norm across channels, and upsample each map to the input
resolution. An image-level score is then obtained by combining the maps with
non-negative weights $\mathbf{w} = (w_1, w_2, w_3)$ on the probability simplex and
reducing the resulting map to a scalar with a pooling operator $\rho$:
%
\begin{equation}
\label{eq:policy}
\begin{aligned}
S(x) &\;=\; \rho\!\left( \textstyle\sum_{k} w_k \, m_k(x) \right), \\
\rho &\in \{\,\text{mean},\ \text{top-}k,\ \max\,\}.
\end{aligned}
\end{equation}
%
We refer to the pair $(\mathbf{w}, \rho)$ as the \emph{scoring policy}. The
published pipeline uses $\mathbf{w} = (\tfrac13,\tfrac13,\tfrac13)$ and
$\rho = \text{mean}$, that is, an unweighted average over all pixels of the sum of
the three maps. This choice is not forced by the method: given a trained detector
the maps $m_k$ are already computed, so any policy can be evaluated on the same
model. Section~\ref{sec:leverage} evaluates the space of policies in
Equation~\eqref{eq:policy} exhaustively, and shows that this reduction is worth
substantially more than the architecture that the search selects for.

\subsection{Protocol and Implementation Details}
\label{subsec:implementation}

Unless stated otherwise, the detector above is trained for 500 epochs with batch
size 64 and seed 42, with a ResNet-18 teacher, on the MVTec-AD
dataset~\cite{bergmann2021mvtec}. The search uses three quick training epochs per
candidate with batch size 16 on the \texttt{wood} category, and is repeated 100
times per run. Inputs are resized to $224\times224$ and normalised with ImageNet
statistics. All figures report image-level AUROC averaged over the 15 categories.
Where a backbone's memory footprint exceeds the device at batch size 64, training
uses gradient accumulation with a smaller micro-batch, keeping the effective batch
size at 64; the micro-batch used is recorded per run and this affects only the
largest backbones.
